

\illustration{0.3}{CarlJacobi.jpg}{
    Carl Gustav Jacob Jacobi (1804--1851), famous brother of Wahnfried Imaginus.
    Public domain on \href{https://commons.wikimedia.org/wiki/File:Carl\_Jacobi.jpg}{Wikimedia Commons}
}





Today, a new paper has been published
in the \emph{Bulletin of Apocryphal Pioneers in Computation}.
According to this paper, the forgotten German number theorist Wahnfried Imaginus Jacobi (1806--1853), while still a secondary student in Potsdam, investigated the decomposition of integers into sums of cubes.
Among the examples noted in the surviving fragments of his notebooks are \[
    2025 = 1^3 + 2^3 + 3^3 + 4^3 + 5^3 + 6^3 + 7^3 + 8^3 + 9^3
\] and the more curious expression \[ 3 = 1^3 + 1^3 + 1^3 = 4^3 + 4^3 + (-5)^3\,,\]
which shows that a solution need not be unique.
Jacobi restricted his attention to small integers and probably did not know the decomposition {\small \[
    3 = 569\,936\,821\,221\,962\,380\,720^3 + (-569\,936\,821\,113\,563\,493\,509)^3 + (-472\,715\,493\,453\,327\,032)^3
\,,\]}%
which was discovered only recently.\footnote{Booker, Andrew R.; Sutherland, Andrew V. (2021), ``On a question of Mordell'', \emph{Proceedings of the National Academy of Sciences}, \textbf{118} (11)}
However, Jacobi did manage to prove that a decomposition into cubes always exists for all positive integers up to $9241$, the $28$th~cuban prime of the first kind.
Although his work was never published, references to the method appear in a marginal annotation in an 1823 letter to his famous brother Carl Gustav Jacob.

\medskip
Given a positive integer~$n$, output a list of at most $10\,000$ integers between $-10\,000$ and $10\,000$
such that the sum of their cubes equals $n$.

